\chapter{Asymptotyka bayesowska: koncentracja posterioru i twierdzenie Bernsteina--von Misesa}
\label{ch:bayesian-asymptotics}
\paragraph{Powiązania i odsyłacze.} Warunki klasycznej asymptotyki zebrano w rozdziale~\ref{ch:asymptotic-estimation}; metody posteriorowe i hierarchiczne w rozdziałach~\ref{ch:bayes-foundations} i~\ref{ch:hierarchical-bayes}.


\section{Dwa rodzaje gwarancji i dwa sposoby granicznego rozumowania}
Rozkład a posteriori jest prawdopodobieństwem warunkowym przy
zadanym priorze i modelu. Nie oznacza to automatycznie, że dla
każdej ustalonej wartości parametru jego przedziały wiarygodności
mają deklarowane pokrycie częstościowe. Zbieżność tych dwóch
opisów pojawia się w wielu \emph{regularnych} modelach przy dużej
próbie, ale nie jest uniwersalnym prawem.

W tym rozdziale rozróżnimy: (i) \emph{zgodność} -- koncentrację
posterioru wokół prawdziwego parametru; (ii) \emph{tempo koncentracji}
-- wielkość otoczeń, w których skupia się posterior;
(iii) \emph{kształt} posterioru po odpowiednim przeskalowaniu.
Dopiero trzeci krok prowadzi do twierdzenia Bernsteina--von Misesa.
Dla przejrzystości ogólne twierdzenie sformułujemy przy jawnych
założeniach lokalnej asymptotycznej normalności i koncentracji
w skali $n^{-1/2}$; wykażemy te warunki w modelach normalnym
i Bernoulliego. Nie będziemy przedstawiać lokalnej normalności
jako wystarczającej bez kontroli ogonów posterioru.

\begin{definicja}[Zgodność posterioru]
\label{def:posterior-consistency}
Niech $(\Theta,d)$ będzie przestrzenią metryczną,
$\theta_0\in\Theta$ prawdziwym parametrem, a
$X_1,X_2,\ldots$ niezależnymi obserwacjami o wspólnym
prawie $P_{\theta_0}$. Niech $\Pi_n(\cdot\mid X_1,\ldots,X_n)$
stanowi posterior na $\Theta$.
Posterior jest \emph{zgodny w punkcie $\theta_0$ według
prawdopodobieństwa}, jeśli dla każdych $\varepsilon,\eta>0$
\[
 P_{\theta_0}\!\left(
 \Pi_n(\{\theta:d(\theta,\theta_0)>\varepsilon\}\mid X_1,\ldots,X_n)
 >\eta\right)\longrightarrow0.
\]
O \emph{zgodności prawie na pewno} mówimy, gdy dla każdego
$\varepsilon>0$ posteriorowa masa powyższego zbioru
zbiega do zera $P_{\theta_0}$-prawie na pewno.
\end{definicja}

\begin{uwaga}
Prawdziwy parametr nie musi być losowany z prioru.
Symbol $P_{\theta_0}$ oznacza rzeczywiste prawo generowania
danych przy ustalonym $\theta_0$; posterior jest losową
miarą, ponieważ zależy od próby. W modelu źle określonym
może nie istnieć żaden $\theta_0$ dający rzeczywisty rozkład
obserwacji; taki przypadek wymaga odrębnej analizy.
\end{uwaga}

\section{Pierwsze dowody zgodności: Bernoulli i skończona przestrzeń parametrów}
\begin{twierdzenie}[Zgodność posterioru beta--Bernoulli]
\label{thm:beta-posterior-consistency}
Niech $a,b>0$, $p_0\in[0,1]$, $X_1,X_2,\ldots$
będą niezależne o rozkładzie $\operatorname{Bernoulli}(p_0)$
i niech prior $p$ będzie równy $\operatorname{Beta}(a,b)$.
Wówczas posterior $\operatorname{Beta}(a+S_n,b+n-S_n)$,
$S_n=\sum_{i=1}^nX_i$, jest zgodny prawie na pewno
w punkcie $p_0$ dla metryki $|p-q|$ na $[0,1]$.
Twierdzenie obejmuje także $p_0=0$ i $p_0=1$,
chociaż sam prior nie przypisuje tym punktom dodatniej masy.
\end{twierdzenie}
\begin{proof}
Oznaczmy przez $m_n$ i $v_n$ średnią oraz wariancję
posterioru. Z twierdzenia~\ref{thm:beta-binomial}
\[
 m_n=\frac{a+S_n}{n+a+b},\qquad
 v_n=\frac{(a+S_n)(b+n-S_n)}{(n+a+b)^2(n+a+b+1)}
 \le\frac1{4(n+a+b+1)}.
\]
Z mocnego prawa wielkich liczb $S_n/n\to p_0$
$P_{p_0}$-prawie na pewno, więc także $m_n\to p_0$.
Dla dowolnego $\varepsilon>0$, gdy
$|m_n-p_0|<\varepsilon/2$, nierówność Czebyszewa
zastosowana \emph{do rozkładu posterioru} daje
\[
 \Pi_n(|p-p_0|>\varepsilon\mid X_1,\ldots,X_n)
 \le\Pi_n(|p-m_n|>\varepsilon/2\mid X_1,\ldots,X_n)
 \le\frac{4v_n}{\varepsilon^2}\longrightarrow0.
\]
Warunek $|m_n-p_0|<\varepsilon/2$ zachodzi ostatecznie
na zdarzeniu prawdopodobieństwa jeden, a oszacowanie
wariancji jest deterministyczne. Dowód nie wykorzystuje
warunku $p_0\in(0,1)$.
\end{proof}

\begin{twierdzenie}[Zgodność przy skończonej liczbie hipotez]
\label{thm:finite-bayes-consistency}
Niech $\Theta=\{\theta_1,\ldots,\theta_K\}$,
$K\ge2$, z priorem $\pi(\theta)>0$ dla każdego $\theta$.
Niech $\theta_0\in\Theta$ będzie prawdziwym parametrem,
$X_i\iid P_{\theta_0}$, a każde $P_\theta$ posiada
gęstość $f_\theta$ względem wspólnej miary $\sigma$-skończonej.
Załóżmy, że dla $\theta\ne\theta_0$ iloraz
$f_\theta(X_1)/f_{\theta_0}(X_1)$ jest dodatni
$P_{\theta_0}$-prawie na pewno, jego logarytm jest
całkowalny oraz
\[
 D(\theta_0\Vert\theta)
 :=\E_{\theta_0}\!\left[
  \log\frac{f_{\theta_0}(X_1)}{f_\theta(X_1)}\right]>0.
\]
Wtedy $\Pi_n(\{\theta_0\}\mid X_1,\ldots,X_n)\to1$
$P_{\theta_0}$-prawie na pewno.
\end{twierdzenie}
\begin{proof}
Dla dowolnego $\theta\ne\theta_0$ posteriorowy iloraz
prawdopodobieństw ma postać
\[
 \frac{\Pi_n(\{\theta\}\mid X)}{\Pi_n(\{\theta_0\}\mid X)}
 =\frac{\pi(\theta)}{\pi(\theta_0)}
 \prod_{i=1}^n\frac{f_\theta(X_i)}{f_{\theta_0}(X_i)}.
\]
Logarytm ilorazu podzielony przez $n$ zbiega
prawie na pewno do $-D(\theta_0\Vert\theta)<0$
na mocy mocnego prawa wielkich liczb. Każdy z $K-1$
ilorazów dąży więc do zera. Sumując te ilorazy,
otrzymujemy
$\Pi_n(\{\theta_0\}\mid X)
 =(1+\sum_{\theta\ne\theta_0}\Pi_n(\{\theta\}\mid X)/
 \Pi_n(\{\theta_0\}\mid X))^{-1}\to1$.
\end{proof}

\begin{uwaga}[Czego nie dowodzi ostatnie twierdzenie]
Skończoność przestrzeni parametrów pozwala przejść od
zbieżności każdego ilorazu do zbieżności ich sumy.
Przy nieskończonej przestrzeni wymagana jest dodatkowa
kontrola złożoności modelu i zachowania prioru. Zgodność
w modelu skończonym nie stanowi dowodu ogólnego twierdzenia
o zgodności dowolnych nieparametrycznych posteriorów.
\end{uwaga}

\section{Dokładny model normalny: co dzieje się z priorem?}
\begin{twierdzenie}[Normalny posterior i asymptotyczna normalność]
\label{thm:normal-conjugate-bvm}
Niech $\sigma^2,\tau^2>0$, $m_0\in\R$ i
$X_1,\ldots,X_n\iid\Normal(\theta_0,\sigma^2)$
przy ustalonym $\theta_0\in\R$.
W analizie bayesowskiej przyjmijmy właściwy prior
$\theta\sim\Normal(m_0,\tau^2)$. Posterior jest normalny
z parametrami
\[
 v_n=(\tau^{-2}+n\sigma^{-2})^{-1},\qquad
 m_n=v_n\left(m_0/\tau^2+n\overline X_n/\sigma^2\right).
\]
Przy $n\to\infty$ zachodzi
$\sqrt n(m_n-\overline X_n)\to0$ w prawdopodobieństwie,
$n v_n\to\sigma^2$ i
\[
 \Law\big(\sqrt n(\theta-\overline X_n)\mid X_1,\ldots,X_n\big)
 \longrightarrow\Normal(0,\sigma^2)
\]
w odległości całkowitej wariacji, w prawdopodobieństwie
$P_{\theta_0}$. Ponadto centralny posteriorowy przedział
wiarygodności poziomu $1-\alpha$ ma pokrycie częstościowe
dążące do $1-\alpha$ dla każdego $\alpha\in(0,1)$.
\end{twierdzenie}
\begin{proof}
Pierwsze wzory pochodzą z uzupełnienia kwadratu
w iloczynie normalnej wiarygodności i prioru.
Zachodzi dokładnie
$m_n-\overline X_n=\sigma^2(m_0-\overline X_n)/
(\sigma^2+n\tau^2)$.
Ponieważ $\overline X_n\to\theta_0$ w prawdopodobieństwie,
mnożenie przez $\sqrt n$ daje granicę zero.
Z kolei $nv_n=\sigma^2n\tau^2/(\sigma^2+n\tau^2)\to\sigma^2$.
Rozkład przeskalowanego parametru jest dokładnie
$\Normal(\sqrt n(m_n-\overline X_n),nv_n)$.
Ciągłość gęstości normalnej względem parametrów
w $L^1(\lambda_1)$, którą otrzymujemy z lematu Scheffégo,
daje zbieżność w całkowitej wariacji.
Centralny przedział ma końce
$m_n\pm z_{1-\alpha/2}\sqrt{v_n}$,
gdzie $z_{1-\alpha/2}$ jest kwantylem standardowego
rozkładu normalnego. Z centralnego twierdzenia granicznego,
$\sqrt n(m_n-\theta_0)/\sigma\Rightarrow\Normal(0,1)$
oraz $\sqrt{nv_n}/\sigma\to1$; zatem prawdopodobieństwo
zdarzenia $|m_n-\theta_0|\le z_{1-\alpha/2}\sqrt{v_n}$
dąży do $1-\alpha$.
\end{proof}

\section{Lokalna asymptotyczna normalność: precyzyjne warunki}
Aby oddzielić własność \emph{wiarygodności} od własności
\emph{posterioru}, wprowadzamy formalną lokalną aproksymację
log-wiarygodności. Jest to warunek, a nie uniwersalna
własność każdego modelu statystycznego.

\begin{definicja}[Lokalna asymptotyczna normalność (LAN)]
\label{def:lan}
Niech $\Theta\subseteq\R^d$ będzie zbiorem otwartym,
$\theta_0\in\Theta$, $d\ge1$, a $\ell_n(\theta)$
logarytmem wiarygodności z próby $n$ obserwacji
w eksperymencie o prawie $P_{\theta_0}^{(n)}$.
Mówimy, że model spełnia warunek LAN w $\theta_0$
z macierzą informacji $I_0\in\R^{d\times d}$ dodatnio
określoną, jeśli istnieją wektory losowe $\Delta_n\in\R^d$
takie, że $\Delta_n\Rightarrow\Normal_d(0,I_0)$ pod
$P_{\theta_0}^{(n)}$ oraz dla każdego $M<\infty$
\[
 \sup_{\norm{h}\le M}\left|
 \ell_n(\theta_0+h/\sqrt n)-\ell_n(\theta_0)
 -h^\top\Delta_n+\frac12h^\top I_0h\right|
 \xrightarrow{P_{\theta_0}}0,
\]
przy czym dla dostatecznie dużego $n$ wszystkie
$\theta_0+h/\sqrt n$ z $\norm h\le M$ należą do $\Theta$.
\end{definicja}

\begin{definicja}[Koncentracja posterioru w skali $n^{-1/2}$]
\label{def:root-n-tightness}
Przy $\Theta\subseteq\R^d$ i prawdziwym $\theta_0$
niech $H_n=\sqrt n(\theta-\theta_0)$ będzie zmienną
posteriorową. Posterior jest \emph{ciasny w skali $n^{-1/2}$
w prawdopodobieństwie}, jeśli dla każdego $\eta>0$
\[
 \lim_{M\to\infty}\limsup_{n\to\infty}
 P_{\theta_0}^{(n)}\!\left(
 \Pi_n(\norm{H_n}>M\mid X^{(n)})>\eta\right)=0.
\]
To warunek mocniejszy niż zgodność w każdym ustalonym
otoczeniu $\theta_0$.
\end{definicja}

\begin{twierdzenie}[Bernstein--von Mises przy jawnych warunkach]
\label{thm:bvm-under-lan}
Niech $\Theta\subseteq\R^d$ będzie otwarty,
$\theta_0\in\Theta$, a dane $X^{(n)}$ pochodziły z
modelu zdominowanego z log-wiarygodnością $\ell_n$.
Zakładamy, że:
\begin{enumerate}
\item model spełnia LAN z macierzą $I_0$ w sensie
definicji~\ref{def:lan};
\item prior ma gęstość $\pi$ względem $\lambda_d$,
ciągłą i ściśle dodatnią w punkcie $\theta_0$;
\item posterior jest właściwy oraz spełnia warunek
z definicji~\ref{def:root-n-tightness}.
\end{enumerate}
Niech $\mathcal Q_n$ oznacza posteriorowy rozkład
$H_n=\sqrt n(\theta-\theta_0)$, a $\mathcal G_n$
rozkład $\Normal_d(I_0^{-1}\Delta_n,I_0^{-1})$.
Wtedy
\[
 \norm{\mathcal Q_n-\mathcal G_n}_{\rm TV}
 :=\sup_{A\in\calB(\R^d)}|\mathcal Q_n(A)-\mathcal G_n(A)|
 \xrightarrow{P_{\theta_0}}0.
\]
Jeżeli ponadto estymator $\widehat\theta_n$ spełnia
$\sqrt n(\widehat\theta_n-\theta_0)=I_0^{-1}\Delta_n+o_{P_{\theta_0}}(1)$,
to posteriorowy rozkład
$\sqrt n(\theta-\widehat\theta_n)$ zbiega w całkowitej
wariacji, w prawdopodobieństwie, do $\Normal_d(0,I_0^{-1})$.
\end{twierdzenie}
\begin{proof}
Dla $h$ takiego, że $\theta_0+h/\sqrt n\in\Theta$,
posterior ma względem $\lambda_d$ gęstość proporcjonalną do
\[
 g_n(h)=\exp\big(\ell_n(\theta_0+h/\sqrt n)-\ell_n(\theta_0)\big)
 \pi(\theta_0+h/\sqrt n).
\]
Czynnik $n^{-d/2}$ po zmianie zmiennych jest wspólny
i skraca się przy normalizacji.
Z warunku LAN i ciągłości prioru, na każdej kuli
$K_M=\{h:\norm h\le M\}$ funkcja $g_n(h)$
po podzieleniu przez stałą $\pi(\theta_0)$ jest
jednostajnie, w prawdopodobieństwie, bliska
$g_n^0(h)=\exp(h^\top\Delta_n-h^\top I_0h/2)$.
Aby sformułować to bez niejawnego ograniczenia $\Delta_n$,
ustalamy $C>0$. Na zdarzeniu $\norm{\Delta_n}\le C$
obie funkcje są na $K_M$ ograniczone z góry i z dołu
przez dodatnie stałe zależne od $C,M,I_0$.
Ponieważ $\Delta_n=O_{P_{\theta_0}}(1)$,
można dobrać $C$ tak, aby dopełnienie tego zdarzenia
miało dowolnie małe prawdopodobieństwo.
Stąd unormowane ograniczenia $\mathcal Q_n$ i $\mathcal G_n$
do $K_M$ (po podzieleniu przez ich masy na $K_M$)
zbliżają się do siebie w całkowitej wariacji.
Istotnie, zbieżność jednostajna ich nieunormowanych
gęstości na $K_M$ pociąga zbieżność całek normalizujących
i następnie gęstości unormowanych w $L^1(\lambda_d)$.
Dla dowolnych miar prawdopodobieństwa $Q,G$ oraz zbioru
$K$ o dodatnich masach mamy oszacowanie
\[
 \norm{Q-G}_{\rm TV}
 \le \norm{Q(\cdot\mid K)-G(\cdot\mid K)}_{\rm TV}
       +Q(K^c)+G(K^c).
\]
Wybieramy $M$ tak duże, aby warunek koncentracji posterioru
zapewniał małą masę $\mathcal Q_n(K_M^c)$ z dużym
prawdopodobieństwem. Jednocześnie
$\mathcal G_n(K_M^c)$ jest mała z dużym prawdopodobieństwem,
ponieważ $\Delta_n=O_P(1)$ i macierz kowariancji
$I_0^{-1}$ jest stała. Następnie stosujemy zbieżność
na kuli, a na końcu pozwalamy zwiększać $M$.
Daje to pierwszą tezę.
Druga wynika przez przesunięcie obu miar o
$-\sqrt n(\widehat\theta_n-\theta_0)$ i z ciągłości
w całkowitej wariacji rodziny normalnej względem wektora
średnich, gdy różnica średnich jest $o_P(1)$.
\end{proof}

\begin{uwaga}[Dlaczego założenie o ogonach jest istotne?]
LAN opisuje wiarygodność tylko dla lokalnych zmian parametru
rzędu $n^{-1/2}$. Nie wyklucza odległych obszarów przestrzeni
parametrów, które zachowują znaczną masę posteriorową.
W zastosowaniach do ogólnych modeli warunek ciasności
w skali $n^{-1/2}$ wyprowadza się zwykle z identyfikowalności,
regularności i istnienia odpowiednich testów lub oszacowań
wiarygodności. Tego głębokiego kroku nie zastępujemy
samym rozwinięciem Taylora.
\end{uwaga}

\section{Model Bernoulliego: pełna weryfikacja warunków LAN i koncentracji}
\begin{twierdzenie}[Bernstein--von Mises dla prioru beta we wnętrzu]
\label{thm:bernoulli-bvm}
Niech $p_0\in(0,1)$, $a,b>0$,
$X_1,\ldots,X_n\iid\operatorname{Bernoulli}(p_0)$,
a prior $p$ ma rozkład $\operatorname{Beta}(a,b)$.
Niech $\widehat p_n=S_n/n$. Wówczas posteriorowy rozkład
$\sqrt n(p-\widehat p_n)$ zbiega w całkowitej wariacji,
w prawdopodobieństwie $P_{p_0}$, do
$\Normal(0,p_0(1-p_0))$.
W szczególności dla $\alpha\in(0,1)$ centralny
posteriorowy przedział wiarygodności poziomu $1-\alpha$
ma asymptotyczne pokrycie częstościowe $1-\alpha$.
\end{twierdzenie}
\begin{proof}
Dla $\ell_n(p)=S_n\log p+(n-S_n)\log(1-p)$
na $(0,1)$ wynik różniczkowania w $p_0$ daje
\[
 \Delta_n=\frac{S_n-np_0}{\sqrt n\,p_0(1-p_0)},
 \qquad I_0=\frac1{p_0(1-p_0)}.
\]
Z centralnego twierdzenia granicznego
$\Delta_n\Rightarrow\Normal(0,I_0)$.
Rozwinięcie Taylora $\ell_n(p_0+h/\sqrt n)$ wokół $p_0$
ma drugi współczynnik
$\ell_n''(p)/n=-S_n/(np^2)-(1-S_n/n)/(1-p)^2$.
Na każdym zwartym otoczeniu $p_0$ wyrażenie to zbiega
jednostajnie w prawdopodobieństwie do
$-p_0/p^2-(1-p_0)/(1-p)^2$, a w $p_0$ przyjmuje
wartość $-I_0$. Dla ograniczonych $h$ otrzymujemy LAN.
Gęstość beta jest ciągła i dodatnia w $p_0$.

Pozostaje sprawdzić \emph{tempo} koncentracji.
Niech $m_n=(a+S_n)/(n+a+b)$ i
$v_n=\Var(p\mid X)\le1/[4(n+a+b+1)]$.
Ponieważ $m_n-\widehat p_n=O(1/n)$ jednostajnie
w $S_n\in\{0,\ldots,n\}$, a
$\sqrt n(\widehat p_n-p_0)=O_{P_{p_0}}(1)$,
mamy $\sqrt n(m_n-p_0)=O_{P_{p_0}}(1)$.
Nierówność Czebyszewa pod posteriorem daje dla $M>0$
\[
 \Pi_n(\sqrt n|p-p_0|>M\mid X)
 \le\frac{4nv_n}{M^2}
 \quad\text{na zdarzeniu }\sqrt n|m_n-p_0|\le M/2.
\]
Ponieważ $nv_n\le1/4$, wybierając najpierw $M$ duże,
a potem $n$ duże otrzymujemy ciasność posterioru
w skali $n^{-1/2}$. Możemy zatem zastosować
twierdzenie~\ref{thm:bvm-under-lan}, a z równości
$I_0^{-1}\Delta_n=\sqrt n(\widehat p_n-p_0)$
wynika pierwsza teza.
Zbieżność całkowitej wariacji do normalnego rozkładu
o dodatniej gęstości implikuje zbieżność posteriorowych
kwantyli po standaryzacji. Z CLT dla $\widehat p_n$
wynika wtedy asymptotyczne pokrycie centralnego
przedziału.
\end{proof}

\section{Co się dzieje na brzegu? Ścisły kontrprzykład normalności}
\begin{twierdzenie}[Posterior przy prawdziwym $p_0=0$]
\label{thm:boundary-beta-gamma}
Niech $a,b>0$ i $X_1,\ldots,X_n\iid\operatorname{Bernoulli}(0)$.
W analizie użyjmy prioru $\operatorname{Beta}(a,b)$,
a więc posterioru $p\mid X\sim\operatorname{Beta}(a,b+n)$.
Wtedy
\[
 \Law(np\mid X_1,\ldots,X_n)
 \Rightarrow\operatorname{Gamma}(a,1),
\]
gdzie $\operatorname{Gamma}(a,1)$ oznacza rozkład gamma
z parametrem kształtu $a$ i intensywnością $1$.
Zatem właściwą skalą w tym modelu brzegowym jest $n^{-1}$,
a granica nie jest normalna.
\end{twierdzenie}
\begin{proof}
Przy $p_0=0$ wszystkie obserwacje są zerowe z
prawdopodobieństwem jeden. Gęstość $U_n=np$ względem
$\lambda_1$ jest dla $0<u<n$ równa
\[
 f_n(u)=\frac{u^{a-1}(1-u/n)^{n+b-1}}{n^a B(a,b+n)}.
\]
Z asymptotyki funkcji gamma, wynikającej np. ze wzoru
Stirlinga, $n^a B(a,b+n)\to\Gamma(a)$.
Dla ustalonego $u>0$ mamy
$(1-u/n)^{n+b-1}\to e^{-u}$, stąd
$f_n(u)\to u^{a-1}e^{-u}/\Gamma(a)$.
Granica jest gęstością probabilistyczną na $(0,\infty)$.
Lemat Scheffégo (zbieżność prawie wszędzie nieujemnych
gęstości o całce jeden do gęstości o całce jeden)
daje nawet zbieżność w $L^1(\lambda_1)$, a zatem
zbieżność w całkowitej wariacji.
\end{proof}

\begin{przyklad}[Próba bez wad nie uzasadnia normalnego przedziału]
Przy $a=b=1$, jeśli nie stwierdzono żadnej wady
w $n$ niezależnych próbach, posterior to
$\operatorname{Beta}(1,n+1)$.
Jego średnia wynosi $1/(n+2)$, natomiast asymptotycznie
$n p$ ma rozkład wykładniczy o średniej jeden.
Symetryczny przedział otrzymany przez naiwną aproksymację
normalną może wychodzić poza $[0,1]$ i pomijać
asymetrię posterioru. W tej sytuacji należy stosować
dokładne kwantyle beta albo odpowiedni jednostronny
przedział, a nie automatycznie wzór oparty na CLT we wnętrzu.
\end{przyklad}

\section{Jak interpretować przedziały i jak rozpoznać zawodność asymptotyki?}
\begin{stwierdzenie}[Zbieżność pokrycia w regularnym modelu jednowymiarowym]
\label{prop:credible-frequentist-coverage}
W założeniach twierdzenia~\ref{thm:bvm-under-lan} dla $d=1$
załóżmy dodatkowo, że istnieje estymator
$\widehat\theta_n$ spełniający jego ostatni warunek.
Niech $[L_n,U_n]$ będzie centralnym posteriorowym
przedziałem wiarygodności o masie $1-\alpha$,
$\alpha\in(0,1)$, gdzie kwantyle są brane względem
posterioru $\theta$.
Wtedy $P_{\theta_0}(\theta_0\in[L_n,U_n])\to1-\alpha$.
\end{stwierdzenie}
\begin{proof}
Po standaryzacji $\sqrt n(\theta-\widehat\theta_n)$
posterior zbiega w całkowitej wariacji do
$\Normal(0,I_0^{-1})$ w prawdopodobieństwie.
Ponieważ dystrybuanta tego rozkładu jest ciągła i ściśle
rosnąca, dolny i górny posteriorowy kwantyl mają postać
$\sqrt n(L_n-\widehat\theta_n)
 =-z_{1-\alpha/2}/\sqrt{I_0}+o_P(1)$ i
$\sqrt n(U_n-\widehat\theta_n)
 = z_{1-\alpha/2}/\sqrt{I_0}+o_P(1)$.
Z założenia dotyczącego estymatora i z LAN
$\sqrt n(\widehat\theta_n-\theta_0)
 \Rightarrow\Normal(0,I_0^{-1})$.
Wobec tego zdarzenie pokrycia jest asymptotycznie
równoważne nierówności
$|\sqrt{nI_0}(\widehat\theta_n-\theta_0)|
 \le z_{1-\alpha/2}+o_P(1)$, której prawdopodobieństwo
zbiega do $1-\alpha$.
\end{proof}

\begin{center}\small
\begin{tabular}{>{\raggedright\arraybackslash}p{.26\textwidth}>{\raggedright\arraybackslash}p{.32\textwidth}>{\raggedright\arraybackslash}p{.27\textwidth}}
\toprule
\textbf{Sytuacja} & \textbf{Czy normalność jest uzasadniona?} & \textbf{Lepszy krok praktyczny}\\
\midrule
Duża próba, parametr we wnętrzu regularnego modelu & Często tak, przy spełnieniu LAN i koncentracji & Porównać kwantyle posterioru z aproksymacją normalną.\\
Prawdopodobieństwo bliskie $0$ lub $1$ & Aproksymacja może być zła, na brzegu bywa fałszywa & Dokładny posterior beta, przedział jednostronny.\\
Model mieszaninowy z zamiennymi etykietami & Lokalna LAN nie wyklucza wielomodalnego posterioru & Parametry identyfikowalne lub funkcjonały niezmiennicze.\\
Mało grup, nieznana wariancja międzygrupowa & Hyperposterior może być bardzo szeroki & Pełny Bayes i analiza wrażliwości hiperprioru.\\
Liczba parametrów rośnie z $n$ & Klasyczne twierdzenie o ustalonym $d$ nie wystarcza & Osobne warunki wysokowymiarowe.\\
Źle określony model & Koncentracja może zajść wokół pseudo-prawdy & Diagnostyka i odporna niepewność.\\
\bottomrule
\end{tabular}
\end{center}

\section{Laboratorium: porównanie skończonej próby z granicą}
W programie porównujemy dokładny posterior beta z rozkładem
normalnym wynikającym z teorii asymptotycznej. Dla małych
$n$ i rzadkich zdarzeń rozbieżności bywają widoczne.
Nie wolno na podstawie samej symulacji wnioskować o
spełnieniu warunków twierdzenia w dowolnym modelu.

\Needspace{23\baselineskip}
\begin{lstlisting}[language=Python,caption={Kompletny program: zbieżność posterioru beta do aproksymacji normalnej i przykład brzegowy (Python, SciPy).}]
import numpy as np
from scipy.stats import beta, gamma, norm

rng = np.random.default_rng(2026)
p0, a, b = 0.07, 2.0, 18.0
z = norm.ppf(0.975)
for n in [25, 100, 1000, 10000]:
    s = rng.binomial(n, p0)
    phat = s / n
    post_a, post_b = a + s, b + n - s
    exact_ci = beta.ppf([0.025, 0.975], post_a, post_b)
    # Operational Wald approximation: plug in the observed proportion.
    gaussian_ci = phat + np.array([-1, 1]) * z * np.sqrt(phat*(1-phat)/n)
    print("n", n, "s", s, "exact", exact_ci, "Wald", gaussian_ci)

# At the boundary p0=0, n*p converges to Gamma(a, rate=1).
n = 10000
post_a, post_b = a, b + n
print("n * posterior median:", n * beta.ppf(0.5, post_a, post_b))
print("limiting Gamma median:", gamma.ppf(0.5, a, scale=1.0))
\end{lstlisting}

\Needspace{23\baselineskip}
\begin{lstlisting}[language=Julia,caption={Kompletny program: posterior normalny i asymptotyczna niepewność średniej (Julia, Distributions.jl).}]
using Random, Statistics, Distributions

Random.seed!(2026)
sigma, tau, m0, theta0 = 2.0, 3.0, 0.0, 1.2
for n in (20, 200, 2000)
    x = rand(Normal(theta0, sigma), n)
    xbar = mean(x)
    vpost = inv(inv(tau^2) + n / sigma^2)
    mpost = vpost * (m0 / tau^2 + n * xbar / sigma^2)
    exact = quantile(Normal(mpost, sqrt(vpost)), [0.025, 0.975])
    asymptotic = xbar .+ [-1.0, 1.0] .* (1.9599639845 * sigma / sqrt(n))
    println((n=n, posterior=exact, asymptotic=asymptotic))
end
\end{lstlisting}

\begin{uwaga}[Aproksymacja operacyjna a znany tylko w symulacji parametr]
W listingu powyżej czynnik wariancji w przedziale Walda obliczamy
z $\widehat p_n$, a nie z $p_0$. W rzeczywistej analizie
$p_0$ jest nieznany. Przy $S_n=0$ operacyjny przedział
Walda degeneruje się do $[0,0]$, podczas gdy dokładny
posterior beta pozostaje nietrywialny. Znajomość $p_0$
w kodzie symulacyjnym służy jedynie generowaniu danych
i ocenie błędu, nie konstruowaniu przedziału.
\end{uwaga}

\begin{zadanie}[Asymptotyka przy nieznanej jakości]
W modelu Bernoulliego z priorem $\operatorname{Beta}(1,1)$
zaobserwowano $S_n$ wad na $n$ niezależnych prób.
Zakładając $p_0\in(0,1)$ i $S_n/n\to p_0$,
obliczyć granicę $n\Var(p\mid S_n)$.
\par\emph{Rozwiązanie.} Z wzoru na wariancję posterioru
\[
 n\Var(p\mid S_n)=
 \frac{n(S_n+1)(n-S_n+1)}{(n+2)^2(n+3)}.
\]
Dzieląc licznik i mianownik przez $n^3$ otrzymujemy
granicę $p_0(1-p_0)$.
To zgadza się z odwrotnością jednostkowej informacji
Fishera $I_0^{-1}=p_0(1-p_0)$; samo takie dopasowanie
wariancji nie dowodzi jeszcze zbieżności całego rozkładu
posterioru do normalnego.
\end{zadanie}

